\chapter{Checklist de reprodutibilidade}
\label{app:checklist-reprodutibilidade}

Este checklist orienta a repetição de uma análise. A reprodução não consiste apenas em obter o mesmo número. Também é necessário usar as mesmas fontes, unidades, regras de seleção e fronteiras. Ao final, outra pessoa deve conseguir explicar de onde veio cada parcela do resultado.

\section{Checklist técnico}
\label{sec:app-checklist-tecnico}

O primeiro passo é localizar os arquivos que definem o documento e a metodologia. A Tabela~\ref{tab:app-checklist-tecnico} indica onde cada parte fica registrada. Antes de executar uma análise, deve-se confirmar que o código, os dados processados e o relatório correspondem ao mesmo conjunto de regras.

\begin{table}[htbp]
\centering
\small
\caption{Checklist técnico de reprodutibilidade do relatório.}
\label{tab:app-checklist-tecnico}
\begin{tabularx}{\textwidth}{L{0.25\textwidth}L{0.32\textwidth}Y}
\toprule
Item & Onde conferir & O que deve ser verificado \\
\midrule
Documento principal & \texttt{docs/report/main.tex} & Ordem dos capítulos, apêndices e base bibliográfica. \\
Capítulos modulares & \texttt{docs/report/chapters/} & Texto da metodologia, implementação, resultados e limites. \\
Apêndices de suporte & \texttt{docs/report/appendices/} & Parâmetros, auditorias e checklists usados para explicar a execução. \\
Figuras e diagramas & \texttt{docs/report/images/} & Correspondência entre a imagem, a legenda e a regra descrita no texto. \\
Tabelas e matrizes & Ambientes \texttt{tabular}, \texttt{tabularx} e \texttt{longtable} & Unidade, fonte e significado de cada indicador quantitativo. \\
Base bibliográfica & \texttt{docs/references.bib} & Existência das referências citadas e compatibilidade com a afirmação. \\
\bottomrule
\end{tabularx}
\end{table}

\section{Sequência recomendada de reprodução}

Uma execução pode ser conferida em seis passos:

\begin{enumerate}
    \item Fixar origem, destino e massa da carga.
    \item Registrar os portos escolhidos e as distâncias rodoviárias.
    \item Listar as viagens ANTAQ em que a origem marítima aparece antes do destino marítimo.
    \item Conferir carga a bordo, distância, trabalho de transporte e fonte da intensidade em cada viagem.
    \item Recalcular a mediana ponderada e, separadamente, conferir o caminho geométrico selecionado.
    \item Somar as pernas representadas e verificar se a saída continua dentro das fronteiras \ttw{} e de custo modelado.
\end{enumerate}

Em Santos--Manaus, essa sequência deve recuperar 89 observações em 22 corredores, sendo 40 por IMO e 49 por fallback. A participação do fallback deve ser 59,54\% do trabalho de transporte. A intensidade aplicada deve ser 9,322050~g/(t$\cdot$nm), e a média diagnóstica, 25,243618~g/(t$\cdot$nm). O caminho deve ser a ligação direta de 6.112~km ou 3.300,216~nm.

Se a massa de teste for 14~t, a verificação manual da parcela de navegação deve produzir:

\begin{equation}
9{,}322050\times14\times3.300{,}216/1000
= 430{,}71~\text{kg de combustível}.
\end{equation}

Somente a massa de 14~t é uma entrada ilustrativa. Os indicadores do par e o caminho são dados da execução atual. O valor de 430,71~kg não inclui toda a cadeia porta a porta.

\section{Checklist metodológico}
\label{sec:app-checklist-metodologico}

A Tabela~\ref{tab:app-checklist-metodologico} deve ser usada como uma lista de confirmação. Cada linha responde ``sim'' ou ``não'' para uma condição necessária. Se uma condição não for atendida, o resultado precisa ser corrigido, classificado como limitação ou excluído da comparação quantitativa.

\begingroup
\footnotesize
\setlength{\LTleft}{0pt}
\setlength{\LTright}{0pt}
\setlength{\tabcolsep}{4pt}
\renewcommand{\arraystretch}{1.10}
\begin{longtable}{L{0.29\textwidth}L{0.65\textwidth}}
\caption{Checklist de consistência metodológica do framework.}
\label{tab:app-checklist-metodologico}\\
\toprule
Dimensão de controle & Pergunta de verificação \\
\midrule
\endfirsthead
\toprule
Dimensão de controle & Pergunta de verificação \\
\midrule
\endhead
Comparação porta a porta & A origem, o destino e a massa são iguais na rodovia direta e na cadeia multimodal? \\
Fronteira ambiental & O resultado está identificado como emissão operacional \ttw{} de \cooe{}, sem ser confundido com \wtw{} ou LCA? \\
Fronteira econômica & O valor está identificado como custo operacional modelado, e não como frete ou tarifa? \\
Sequência das viagens ANTAQ & Todas as escalas pertencem à mesma viagem, com a origem antes do destino? \\
Cobertura de corredores & Ligações diretas e multiescala foram consideradas, com uma contribuição por viagem e par? \\
Carga a bordo & Embarques e desembarques foram acumulados e a carga permaneceu não negativa? \\
Trabalho de transporte & Foram somados carga a bordo vezes distância de todos os subtrechos do OD? \\
Intensidade marítima & O sistema procurou primeiro o IMO no EU MRV e registrou qualquer fallback por classe ou tipo? \\
Outliers do fallback & O tipo usou média aparada de 1\% ou mediana; a classe usou média aparada ou mediana; e o IMO exato foi preservado? \\
Agregação do par & A intensidade é a mediana inferior ponderada pelo trabalho de transporte de todas as viagens e corredores elegíveis? \\
Pesos nulos & Pesos nulos foram excluídos quando havia peso positivo e tratados por mediana não ponderada quando todos eram nulos? \\
Seleção de caminho & Foi escolhida uma ligação direta observada ou, na sua ausência, a menor cadeia completa, sem combinar viagens? \\
Proveniência de distâncias & Cada perna informa se veio de matriz, API, referência, fallback ou override? \\
Componentes portuários & As operações portuárias têm fonte e o \emph{hoteling} não foi duplicado quando se usou intensidade MRV abrangente? \\
Tratamento de dados ausentes & Dados indisponíveis permaneceram como lacuna, em vez de serem transformados em zero? \\
Compatibilidade do artefato & A matriz informa schema, data, política de intensidade, regra de caminho e cobertura? \\
Classificação de evidência & O resultado foi limitado à sua categoria, sem alegação de superioridade universal? \\
\bottomrule
\end{longtable}
\endgroup

Ao concluir o checklist, o leitor deve conseguir repetir a fórmula e também explicar seus limites. No exemplo de 14~t, isso significa reproduzir os 430,71~kg de navegação e declarar que o número não é emissão \wtw{}, custo comercial ou confirmação de serviço.
